\documentclass{article}
\usepackage[main]{neurips_2026}
\renewcommand{\ttdefault}{pcr}
\usepackage[utf8]{inputenc}
\usepackage[T1]{fontenc}
\usepackage{amsmath,amsfonts,amssymb}
\usepackage{graphicx,booktabs,longtable,microtype}
\usepackage{xcolor}
\usepackage{hyperref,url}
\title{Evaluating Joint Reliability for Observed Compositions}
\author{Anonymous Authors}
\newcommand{\simplex}{\Delta^{K-1}}
\newcommand{\R}{\mathbb{R}}
\input{figures/diagnostic_captions}
\input{figures/core_captions}
\input{figures/region_captions}
\input{figures/final_captions}
\begin{document}
\maketitle
\begin{abstract}
Conformal regions for compositions must distinguish joint coverage of an observed vector from coverage of class labels or latent probabilities. We compare split-conformal score geometries, training-only scaling and fixed-group calibration across 19 frozen synthetic cells with 100 independently regenerated/refitted trials each. A common original-simplex area and separate numerical audit make efficiency comparisons interpretable. With ten components, global log-ratio calibration has population coverage .9003 [95\% interval .8975,.9030] but high-variance minority coverage .0511 [.0371,.0652]; fixed-group calibration raises minority coverage to .9096 [.9008,.9184]. At three components, the log-ratio-minus-Euclidean area difference changes from -.0439 in the interior to +.0498 near the boundary, with enlarged paired intervals excluding zero, while TV/Euclidean ordering remains unresolved. Group calibration tolerates the specified prevalence shift but returns the full domain in 31 of 100 scarce-minority calibrations. Nearly full-dimensional area also coexists with about 60\% coverage on zero faces. These conditional findings support evaluating joint events, groups, applicability and numerical resolution together; they establish no universal score winner, latent-probability or real-data guarantee.
\end{abstract}

\input{manuscript/introduction}
\input{manuscript/related-work}
\input{manuscript/evaluation}
\input{manuscript/methods}
\input{manuscript/results}
\input{manuscript/scope}
\clearpage
\bibliographystyle{plainnat}
\bibliography{context/references}
\clearpage
\appendix
\input{manuscript/technical}
\clearpage
\input{manuscript/diagnostics}
\clearpage
\input{manuscript/process}
\clearpage
\input{manuscript/supplementary-results}
\clearpage
\input{manuscript/checklist}
\end{document}
